\chapter{How to Prove That Invisible Things Exist}

\begin{kaichang}
One summer evening, your mother sprays a little floral toilet water in the living room, then goes to the kitchen. You are sitting on the sofa three meters away. Less than half a minute later, a cool fragrance reaches your nose.

Stop and consider how strange this is: the toilet water is still in its bottle on the table, and nobody has brought it over to you, yet you can smell it. \textbf{Something has escaped from the bottle and crossed three meters of air}---but however wide you open your eyes, you cannot see anything.

Does it exist? Of course: your nose is a witness. But a nose cannot speak, produce a photograph, or put the ``fragrance'' on a balance. This leads to one of science's most troublesome and fascinating questions: \textbf{if we cannot see, touch, or photograph something, what entitles us to say it exists?} In this chapter, we learn the art of ``building a case against the invisible.'' Afterward, you will find that the same logic applies to everything from a fragrance to atoms tens of thousands of times smaller than bacteria.
\end{kaichang}

\section{The Visible and the Invisible}

First, let us be clear: the vast majority of things science studies are invisible to the naked eye. Molecules, bacteria, wind, and temperature are invisible; even oxygen cannot be seen, however hard you stare.

Must science therefore rely on guesswork? Quite the opposite. Science has an extremely practical approach: \textbf{instead of looking at the thing itself, look at the traces it leaves}.

You already know how to do this. When you wake on a winter morning and see moisture on the window, you immediately know ``it was cold outside last night.'' You did not stand outside to feel it, but the moisture is evidence. Coming home from school, the smell of cooking in the hallway tells you someone is making dinner. Moving leaves tell you ``the wind has picked up.'' Nobody throughout history has ever seen wind itself; all we see are things it moves---leaves, flags, and ripples on water.

Scientists distinguish two levels of this reasoning:

\begin{itemize}[itemsep=2pt]
\item \textbf{Direct observation}: your senses receive the phenomenon itself, with simple instruments extending them when needed. Examples include a balance needle moving, test paper changing color, or the liquid column in a thermometer rising.
\item \textbf{Indirect evidence}: the object of study is invisible, but its \emph{effects} can be measured. Nobody has seen an electron, for example, but electrons make a fluorescent screen glow when they hit it. The light is visible, and in this sense the electrons have been ``seen.''
\end{itemize}

The crucial point is that one piece of indirect evidence may be unreliable: another explanation could fit. Moisture on the glass might mean the room is too humid; the food smell might come through someone else's extractor fan. The real work of science is therefore \textbf{joining many independent pieces of indirect evidence into a chain}, so that all point to the same explanation and exclude alternatives. The more complete the chain, the more confidently we can say ``the invisible thing exists''---even more confidently than ``I saw it myself,'' because eyes can deceive, while independent measurements are far less likely to deceive together.

\section{The Detective's Six Tools}

To build a case against the invisible, you need suitable tools. These six are the most common in the matter detective's toolbox. You will use them repeatedly throughout the book.

\textbf{First: mass.} A balance is an honest witness that takes no part in arguments. Let mass tell you whether something has arrived or departed. Leave an iron nail in damp air and weigh it again several days later: it has become heavier. Where did the extra mass come from? Something invisible---a component of the air---must have joined the nail. Chapter 1 described ``whether a new substance forms'' as a clue to chemical change, and a gain or loss of mass is often the strongest version of that clue.

\textbf{Second: volume.} Gas is invisible, but it occupies space. Pump air into a balloon, and it expands. Block a syringe's opening and push the plunger: the farther you push, the harder it becomes, because the invisible stuff inside resists. Measuring volume turns ``air occupies space and can be compressed'' into numbers instead of an empty assertion.

\textbf{Third: density.} Divide mass by volume to obtain density, a substance's ``build and physique code.'' For equal-sized pieces, iron is much heavier than aluminum, which is much heavier than wood. Two metal blocks may look identical, but weighing them and measuring their volumes to calculate density immediately reveals different materials inside. Archaeologists use the same yardstick to detect adulterated gold coins.

\textbf{Fourth: temperature.} Temperature tells you how vigorously particles are ``bustling about'' on average (Chapter 4 will explain this fully). Chemical reactions often involve temperature changes: quicklime becomes hot on contact with water, while ammonium nitrate dissolving in a cold pack makes it cool. A thermometer is a probe into the world of particles.

\textbf{Fifth: acidity and alkalinity.} Measure these with pH paper or a pH meter. pH is a number that roughly indicates the concentration of particles called hydrogen ions in a solution. Small numbers (below 7) indicate acidity, large ones (above 7) alkalinity, and values near 7 neutrality (this standard applies only near room temperature; Chapter 32 will explain its limits). Lemon juice, soapy water, stomach fluid, and rainwater each have their own pH; measurement reveals their character.

\textbf{Sixth: spectra.} This is the most remarkable tool and deserves a little more explanation. Send light through a prism or diffraction grating, and it spreads into a rainbow. But heat an element until it glows, or excite it electrically, and its light spreads into a few bright lines at particular positions and colors instead of a continuous rainbow. Each element has its own line pattern, like a fingerprint. The red of neon lights and the yellow of sodium lamps each come from a line in this ``fingerprint.'' In 1868, astronomers found bright lines in the Sun's spectrum that matched no known element. They announced an element on the Sun not yet seen on Earth and named it ``helium,'' from the Greek word for the Sun. Only twenty-seven years later was helium found on Earth. \textbf{An element was ``seen'' by spectroscopy more than ninety million kilometers away before it was found on Earth.} Why spectral lines form fingerprints lies in the arrangement of electrons within atoms---the story of Chapter 9.

Notice the common pattern: \textbf{none of these six tools ``directly sees'' the object of study; all measure effects it produces.} Being invisible is fine, as long as something is measurable.

\section{Making the Invisible Give Itself Away}

With tools in hand, let us examine one of the finest cases of ``indirect conviction'' in scientific history.

\begin{zhengju}
In 1827, the British botanist Brown sprinkled pollen on water and examined it through a microscope. He found pollen particles jiggling without end: left and right, starting and stopping, with no pattern in their directions. At first he thought the pollen was ``alive,'' but finely ground stone and glass behaved the same way. Any sufficiently small particles in water jiggled like this. If life was not the cause, what was pushing them?

The puzzle remained unresolved for nearly eighty years. Some suggested currents caused by uneven water temperature, others illumination. Experiments ruled these explanations out: at constant temperature and in darkness, the particles still moved, and smaller particles and warmer water meant more vigorous motion.

In 1905, Einstein offered a bold quantitative explanation. Water consists of innumerable invisible molecules in constant, jostling motion. A suspended pollen particle is struck from every direction at every moment. For a sufficiently large particle, impacts from different directions cancel, producing no visible motion. But pollen is only a few micrometers across: at one instant, a few more molecules may happen to hit from the left than from the right, giving it a ``shove''; the next instant, the shove comes from elsewhere. The apparently patternless motion is actually \textbf{statistical fluctuations in countless molecular impacts}. Einstein also calculated how far particles should move on average and the precise mathematical relationships with temperature, particle size, and water viscosity.

The theory was beautiful, but nobody had seen its proposed water molecules. Why believe it? Beginning in 1908, the French physicist Perrin carried out a series of painstaking experiments. He made uniformly sized resin particles, tracked hundreds individually under a microscope, and recorded thousands upon thousands of positions. Their actual motion matched Einstein's predictions exactly. Better still, a completely different measurement---counting how suspended particles were distributed with height---let Perrin infer the number of particles per mole, Avogadro's constant. His result, about $6\times 10^{23}$, agreed with other methods.

By then even the most stubborn skeptics conceded. Ostwald, a chemist who had opposed the real existence of atoms and molecules, publicly acknowledged that Brownian motion and Perrin's experiments ``entitled even the most cautious scientists to speak of matter as composed of discrete particles.'' Notice the chain's structure: \textbf{visible pollen jiggling → explained only by water-molecule impacts (after excluding currents, uneven temperature, and other alternatives) → mathematical predictions verified by precise measurements → several independent routes yielding the same number}. Nobody had ``seen'' a water molecule, yet this chain was more reliable than naked eyes. Perrin received the 1926 Nobel Prize in Physics for this work.
\end{zhengju}

The same approach twice rewrote the history of atomic models. One case involved cathode rays: high voltage applied across an evacuated glass tube produces an invisible beam that makes a fluorescent screen glow, bends in a magnetic field, and can turn a tiny paddle wheel. Thomson concluded that the beam was a stream of minute, negatively charged particles---electrons. The other was the gold-foil experiment. Rutherford directed helium nuclei at extremely thin gold foil, expecting all to pass straight through. Instead, a tiny fraction rebounded, like shells fired at paper bouncing back. He inferred that an atom's mass and positive charge were concentrated in an astonishingly small nucleus. Chapter 6 will tell these stories in detail. For now, remember their shared feature: \textbf{the rays and nuclei were invisible, but visible fluorescence, deflection, and rebound established their existence}.

\section{Rules for Good Experiments}

A single observation may be coincidence; one measurement may be a mistake. To turn indirect evidence into a chain of evidence, experiments must follow certain rules. These are lessons learned from centuries of being misled, not scientists' fussiness.

\textbf{Rule one: reproducibility.} An experiment you performed once that nobody else can reproduce proves nothing. In 1989, two chemists announced ``cold fusion'' in an electrochemical cell at room temperature, making headlines worldwide. Dozens of laboratories tried to repeat it, but none succeeded, and the conclusion was withdrawn. Reproducibility is science's first checkpoint for every claim.

\textbf{Rule two: controls.} To discover whether ``fertilizer increases rice yields,'' spreading fertilizer and obtaining a good autumn harvest proves little: perhaps rainfall was especially favorable that year. The right approach is to divide a field into two plots with identical conditions, fertilizing one but not the other. The unfertilized plot is the \textbf{control group}. The single condition that differs is the \textbf{variable}. Ideally, a good experiment \emph{changes only one variable at a time}; otherwise, you cannot tell what deserves credit for the result.

\begin{shiyan}
\textbf{Weighing ``Air That Escaped'': A Controlled Evaporation Experiment}

Materials: two identical small cups (or bowls), tap water, a small piece of plastic wrap, a rubber band, and a digital kitchen scale (a marker will do if you have no scale).

Procedure:
\begin{enumerate}[itemsep=1pt]
\item Put equal amounts of water in both cups (half full is enough).
\item Leave one open; cover the other with plastic wrap secured by a rubber band.
\item Weigh each and record the result (or mark the initial water levels on the cups).
\item Put the cups side by side on a windowsill or elsewhere indoors and leave them for two to three days.
\item Weigh them separately again (or inspect their water levels).
\end{enumerate}

What to observe: How much mass did the open cup lose? What about the covered cup? If the covered cup's mass barely changes while the open cup becomes noticeably lighter, where did the missing mass go?

Safety: This uses only clean water at room temperature and presents no hazards. Put the cups somewhere they will not be knocked over.

Here the covered cup is the control, and ``open or covered'' is the only variable. The open cup's loss of mass forces you to accept that some water became invisible vapor and escaped through the air. You did not see it leave, but the scale did. This is precisely the central skill of this chapter.
\end{shiyan}

\textbf{Rule three: use blinding when needed.} Humans are easily deceived by their expectations. If a doctor knows which patients received real medicine and which received identical-looking starch pills (placebos), that knowledge can bias observations and judgments. Patients who know they have taken a ``new medicine'' also often feel better. Modern drug trials therefore use \textbf{double blinding}: neither patients nor the doctors observing them know who received what. A third party holds the allocation list until the statistical analysis is complete. You may think ``I am not that easily fooled,'' but evidence shows that everyone is subject to expectation effects. Blinding protects honest people.

\textbf{Rule four: correlation is not causation.} This is the easiest rule to stumble over. Statistics show that months with high ice-cream sales also have more drownings. Does ice cream cause drowning? Obviously not. A common cause---hot weather---lies behind both: more people eat ice cream and more people swim in rivers. Sales and drownings are merely \emph{correlated}; no causal arrow connects them. Likewise, roosters crow and the Sun rises together every day, yet you would never think crowing summons the Sun. Establishing causation requires controlled experiments excluding alternatives or a mechanism with an unbroken chain of links. Two curves rising and falling together are nowhere near enough. Remember this rule: you will need it daily when reading news, advertisements, and health-product claims.

\section{Every Measurement Comes with ``Approximately''}

Another admission is necessary: \textbf{no measurement is absolutely exact}. Rulers have a smallest division, scales have a resolution, and your eyes are estimating the last digit. ``This cup contains $200\,\mathrm{mL}$ of water'' honestly means ``approximately $200\,\mathrm{mL}$, within a few milliliters.'' Science calls this margin \textbf{uncertainty}.

The number of digits in a reported measurement carries information. ``$2\,\mathrm{g}$,'' ``$2.0\,\mathrm{g}$,'' and ``$2.00\,\mathrm{g}$'' are three different statements: the first claims precision only to the units digit, the second to 0.1 gram, and the third says your scale can distinguish 0.01 gram. These meaningful digits are called \textbf{significant figures}. Adding a digit exaggerates; omitting one wastes information. Recording only ``2 grams'' when a balance reads to 0.01 gram is using a fine instrument as a crude one.

\begin{qiaoliang}
\textbf{If one object is too small to measure, measure a batch and divide}---an age-old answer to ``too small.'' Let us put numbers to it.

A grain of rice has a mass of about $0.02\,\mathrm{g}$. A household kitchen scale typically resolves $1\,\mathrm{g}$. Put one grain on it, and it reads $0$---not because rice has no mass, but because the grain is fifty times smaller than the scale's resolution. The scale cannot see it.

Try another approach: count 500 grains and weigh them together. The reading is about $10\,\mathrm{g}$. The scale's error is approximately $\pm 1\,\mathrm{g}$, a relative error of about $10\,\%$ for $10\,\mathrm{g}$. Thus each grain has mass $=10\,\mathrm{g}\div 500=0.02\,\mathrm{g}$, with error also reduced to about $10\,\%$. \textbf{If one cannot be measured accurately, measure five hundred and average.}

This simple idea is exactly how scientists handle atoms. One atom has a mass on the order of $10^{-23}$ grams, beyond the reach of any balance. But weigh an astronomically large counted collection and divide by the count, and you obtain the mass of one atom. The problem is how to ``count'' astronomical numbers of atoms. That is the difficulty addressed by Chapter 5's main character, the mole. Perrin's experiment counting Brownian particles was among the earliest successful attempts.
\end{qiaoliang}

\section{Writing Measurements as Symbols}

Numbers need a shared notation; otherwise your ``handspan'' and mine will not balance the same accounts. This is science's symbolic language. Its rules are few but followed worldwide:

\begin{itemize}[itemsep=2pt]
\item \textbf{Attach units to numbers.} ``The nail weighs $5$'' is incomplete---$5$ what? Grams? Kilograms? Write it fully: $m = 5.2\,\mathrm{g}$. Units are not decoration. Wrong units mean a wrong answer; bridges and rockets have suffered accidents for this reason.
\item \textbf{Use scientific notation for very large and very small numbers.} A drop of water contains about $1.7\times 10^{21}$ molecules; a hydrogen atom has a mass of about $1.7\times 10^{-24}\,\mathrm{g}$. The exponent $21$ indicates 21 places after the decimal point; the exponent $-24$ indicates a separation of 24 places before the decimal point. This notation fits astronomical and microscopic numbers on a single line. You will use it constantly in Chapter 5.
\item \textbf{For careful work, state the uncertainty too.} For example, $m = (2.00 \pm 0.05)\,\mathrm{g}$ means the true value is likely between $1.95$ and $2.05$ grams. The ``$\pm$'' is part of the measurement, not modesty. Without it, others cannot judge the strength of your evidence.
\end{itemize}

You may have noticed that our small matter of a fragrance has now been translated into symbolic records with units and error ranges: ``The open cup's mass changed from $152.3\,\mathrm{g}$ to $149.1\,\mathrm{g}$; the covered cup's from $151.8\,\mathrm{g}$ to $151.7\,\mathrm{g}$.'' In this way, invisible things leave testimony on paper.

\section{Even Chains of Evidence Have Limits}

The art of ``looking at effects rather than the thing itself'' is tremendously powerful, but it has clear boundaries. Cross them, and your reasoning fails.

\textbf{First, instrument readings always require interpretation before becoming evidence.} We interpret a rising thermometer column as ``getting warmer,'' but the rise itself is only alcohol expanding. The link between expansion and temperature was established by calibration. More sophisticated instruments involve more layers of interpretation. Before reading an instrument, ask: what does it actually measure?

\textbf{Second, every measurement has a range; conclusions outside it are invalid.} A scale resolving $1\,\mathrm{g}$ reads $0$ for a grain of rice, but its mass is not zero. An ordinary microscope cannot show viruses, but that does not mean they do not exist: they lie below the optical microscope's limit, set by light's wavelength, and became ``visible'' only with electron microscopes. ``Not detected'' and ``nonexistent'' are different things. Half of scientific progress has come from building a more sensitive ``balance.''

\textbf{Third, quantities such as temperature and pH measure collective behavior, not an individual particle's state.} In water at $25\,^{\circ}\mathrm{C}$, some molecules move faster and others slower. Temperature represents the average motion of enormous numbers of molecules. Describing one molecule by temperature is like describing a particular person by the average salary of an entire city: the concept is misplaced.

\textbf{Fourth, even strong indirect evidence remains open to a better explanation.} This is scientific honesty, not weakness. Brownian motion establishes molecular existence very firmly, but science always leaves the door open to new evidence. The more complete the evidence chain, however, the stronger the evidence needed to overturn it.

\begin{wuqu}
\textbf{Misconception: ``Seeing is believing; invisible things do not exist.''} By this standard, air and bacteria do not exist, and the history of infectious disease before bacteria and viruses were discovered becomes inexplicable. In reality, nearly all of science's greatest discoveries concern invisible things: atoms, electrons, genes, and electromagnetic waves. Existence is judged by reproducible, measurable effects that rule out alternatives, never by naked-eye visibility. Beware the opposite extreme too: ``Numbers on an instrument's screen are ultimate truth.'' Instruments have ranges, errors, and layers of interpretation. A solution's own color can interfere with test-paper color changes, and even the most expensive instrument speaks only within its design limits. The detective trusts neither ``I did not see it, so it did not happen'' nor ``the instrument is always right,'' but a chain of evidence that is \textbf{independent, mutually corroborating, and reproducible}.
\end{wuqu}

\section{Back to the Bottle of Fragrance}

Let us reopen the ``missing fragrance'' case from the beginning.

The toilet water contains alcohol and fragrant substances. Their molecules move constantly at the liquid surface, and some overcome their neighbors' pull and escape into the air: evaporation, a process Chapter 4 will examine closely. The escaped molecules wander and collide in the air (the same mechanism behind the pollen motion Brown showed us), spreading throughout the room within minutes. A few fortunate fragrance molecules enter your nose, strike olfactory cells, and trigger a chain of chemical signals to your brain. You ``smell'' them.

Your eyes supplied no evidence at any stage. Yet the chain is complete: mass loss in the open-cup experiment shows that something left the liquid and entered the air; fragrance spreads faster at higher temperatures, consistent with more vigorous molecular motion; move to another room and close the door, and the fragrance cannot reach you, showing its dependence on air movement and molecular diffusion. Every link is measurable and reproducible and excludes all alternative explanations other than ``psychological suggestion.'' \textbf{Invisible, yet supported by overwhelming evidence.}

\section{Where This Skill Is Useful}

You may think indirect evidence and control groups are remote from daily life, but they shape it every day. Every medicine you take must pass randomized, double-blind, placebo-controlled clinical trials before reaching the market. The rules in this chapter distinguish ``seems effective'' from ``actually effective.'' Countless appealing treatments fail this test, and that is a good thing. The air-quality index on your phone comes from continuous measurements of invisible particles. Weather-station temperature records, hospital test reports, and archaeological carbon-14 dating all use the art of ``measure effects, infer the thing itself.'' Even an advertisement claiming ``nine out of ten users approve'' can often be dismantled on the spot by asking ``where is the control group?''

The chapter's deeper purpose is to prepare the way for the book. Nearly everything ahead---atoms, electrons, chemical bonds, molecular machines inside cells---is invisible. You can study them with confidence because you now possess the tools for building the case: \textbf{observe effects, establish controls, calculate errors, and connect the evidence}. Next we use them on a question close at hand: how much of the matter around you is actually ``pure''? Chapter 5 tackles counting atoms, and Chapter 6 shows how the law of definite proportions, cathode rays, and particles rebounding from gold foil gradually \textbf{forced} scientists toward the now-obvious conclusion that atoms exist. It is this art's most glorious case history.

\begin{tiaozhan}
\textbf{Why does averaging repeated measurements improve accuracy?} (Skipping this will not interrupt the main story.)

Suppose each measurement has random error, sometimes high and sometimes low, with no bias on average. What happens when we average $N$ independent measurements? Statistics answers that the typical error of the average decreases approximately as $1/\sqrt{N}$. An error of $\pm 1$ for one measurement becomes about $\pm 1/2$ after averaging 4 and about $\pm 1/10$ after averaging 100. Gaining another digit of precision requires multiplying the measurement count by 100---which is why precision experiments are expensive and slow. By tracking hundreds of particles and recording thousands of positions, Perrin was essentially buying overwhelming precision with a huge $N$. The $1/\sqrt{N}$ rule has a beautiful consequence: large numbers of random events have small relative fluctuations. A mole of gas contains $6\times 10^{23}$ jostling molecules; fluctuations amount to only $\sqrt{N}$, or a fraction on the order of one part in $10^{12}$---too small to measure. Thus the air pressure and water temperature we experience seem rock-steady. \textbf{The stability of the macroscopic world is the statistical outcome of frantic microscopic motion.}
\end{tiaozhan}

\section*{Try It Yourself}

\begin{sikaoti}
\item A lake freezes in winter, but the water beneath the ice remains liquid. What ``invisible witnesses'' can you think of that help fish survive the winter? Give at least two measurable effects. (Hint: think about temperature and density.)
\item An advertisement says: ``City residents who regularly drink this herbal tea have a cold rate $30\,\%$ below the city average.'' List at least two explanations besides ``the tea works,'' and design an experiment to exclude them. State the control group, variables, and whether blinding is needed.
\item Draw a particle diagram of water in an open beaker and the air above it. Use dots for water molecules and show molecules escaping the surface and jostling through the air. Then draw the same beaker with a sealed lid. Add a note: dots and paths are human representations, not what molecules really look like.
\item Estimate a paper clip's mass using a kitchen scale with $1\,\mathrm{g}$ resolution. Design a procedure and calculate the approximate relative error (one paper clip is about $0.5\,\mathrm{g}$). How could you simplify the procedure using an electronic balance with $0.01\,\mathrm{g}$ resolution?
\item Why do smaller particles show more conspicuous Brownian motion? Explain using ``statistical fluctuations in molecular impacts from every direction.'' (Hint: think about $\sqrt{N}$ in the challenge box. With fewer impacts, do different directions cancel more or less completely?)
\item This chapter introduced six detective tools. To show that oxygen in air participates in iron rusting, which would you choose, and what would each measure? Draw your evidence chain as a flowchart: put a measurement result in each box and label each arrow ``from this we infer.''
\end{sikaoti}
